\documentclass[10pt,a4paper]{amsart}
\usepackage{amsmath,amssymb,amsthm,graphicx}
\usepackage{verbatim}
\newenvironment{reference}{\comment}{\endcomment}
%\newenvironment{reference}{}{}
\newenvironment{slogan}{\comment}{\endcomment}
\newenvironment{history}{\comment}{\endcomment}

% For commutative diagrams we use Xy-pic
\usepackage[all]{xy}

% We use 2cell for 2-commutative diagrams.
\xyoption{2cell}
% Theorem environments.
%
\theoremstyle{plain}
\newtheorem{theorem}[subsection]{Theorem}
\newtheorem{proposition}[subsection]{Proposition}
\newtheorem{lemma}[subsection]{Lemma}

\theoremstyle{definition}
\newtheorem{definition}[subsection]{Definition}
\newtheorem{example}[subsection]{Example}
\newtheorem{exercise}[subsection]{Exercise}
\newtheorem{situation}[subsection]{Situation}

\theoremstyle{remark}
\newtheorem{remark}[subsection]{Remark}
\newtheorem{remarks}[subsection]{Remarks}

\numberwithin{equation}{subsection}

% Macros
%
\def\lim{\mathop{\mathrm{lim}}\nolimits}
\def\colim{\mathop{\mathrm{colim}}\nolimits}
\def\Spec{\mathop{\mathrm{Spec}}}
\def\Hom{\mathop{\mathrm{Hom}}\nolimits}
\def\Ext{\mathop{\mathrm{Ext}}\nolimits}
\def\SheafHom{\mathop{\mathcal{H}\!\mathit{om}}\nolimits}
\def\SheafExt{\mathop{\mathcal{E}\!\mathit{xt}}\nolimits}
\def\Sch{\mathit{Sch}}
\def\Mor{\mathop{\mathrm{Mor}}\nolimits}
\def\Ob{\mathop{\mathrm{Ob}}\nolimits}
\def\Sh{\mathop{\mathit{Sh}}\nolimits}
\def\NL{\mathop{N\!L}\nolimits}
\def\CH{\mathop{\mathrm{CH}}\nolimits}
\def\proetale{{pro\text{-}\acute{e}tale}}
\def\etale{{\acute{e}tale}}
\def\QCoh{\mathit{QCoh}}
\def\Ker{\mathop{\mathrm{Ker}}}
\def\Im{\mathop{\mathrm{Im}}}
\def\Coker{\mathop{\mathrm{Coker}}}
\def\Coim{\mathop{\mathrm{Coim}}}

% Boxtimes
%
\DeclareMathSymbol{\boxtimes}{\mathbin}{AMSa}{"02}

%
% Macros for moduli stacks/spaces
%
\def\QCohstack{\mathcal{QC}\!\mathit{oh}}
\def\Cohstack{\mathcal{C}\!\mathit{oh}}
\def\Spacesstack{\mathcal{S}\!\mathit{paces}}
\def\Quotfunctor{\mathrm{Quot}}
\def\Hilbfunctor{\mathrm{Hilb}}
\def\Curvesstack{\mathcal{C}\!\mathit{urves}}
\def\Polarizedstack{\mathcal{P}\!\mathit{olarized}}
\def\Complexesstack{\mathcal{C}\!\mathit{omplexes}}
% \Pic is the operator that assigns to X its picard group, usage \Pic(X)
% \Picardstack_{X/B} denotes the Picard stack of X over B
% \Picardfunctor_{X/B} denotes the Picard functor of X over B
\def\Pic{\mathop{\mathrm{Pic}}\nolimits}
\def\Picardstack{\mathcal{P}\!\mathit{ic}}
\def\Picardfunctor{\mathrm{Pic}}
\def\Deformationcategory{\mathcal{D}\!\mathit{ef}}
\usepackage{cleveref}
\crefname{originallemma}{original lemma}{original lemmas}
\Crefname{originallemma}{Original lemma}{Original lemmas}
\makeatletter
\newlabel{formal-spaces-lemma-base-change-finite-type}{{0AQ5}{source}{}{sourcepointer-1}{}}
\newlabel{formal-spaces-lemma-base-change-finite-type@cref}{{[originallemma][0][]0AQ5}{[1][0][]source}{}{}{}}
\makeatother
\crefname{originalreference}{original reference}{original references}
\Crefname{originalreference}{Original reference}{Original references}
\makeatletter
\newlabel{lemma-flat-and-diagonal-rig-surjective}{{0GD9}{source}{}{sourcepointer-2}{}}
\newlabel{lemma-flat-and-diagonal-rig-surjective@cref}{{[originalreference][0][]0GD9}{[1][0][]source}{}{}{}}
\makeatother
\title{Rig-surjectivity of f with rig-etale rig-surjective diagonal is tested on DVR points}
\author{The Stacks Project Authors}
\date{stacks-project:restricted.tex:local-snapshot-20261012; GNU FDL 1.2 or later}
\begin{document}
\maketitle
\noindent\textit{Editorial note.} All delivered mathematical text is the author's own, copied from the pinned Stacks Project chapter file restricted.tex; no statement or proof step was rewritten, reordered or removed.\par\medskip
\noindent\textit{Editorial note.} The counted proof is the author's complete argument: the proof reduces rig-surjectivity to the flat-locus dichotomy, checks the diagonal hypothesis passes to the base change, and uses the factoring DVR extension to exclude the empty case\par\medskip
\noindent\textit{Editorial note.} The label lemma-rig-monomorphism-rig-surjective names the counted statement, and the proof environment immediately following it in the source is delivered verbatim as its proof.\par\medskip
\noindent\textit{Editorial note.} The counted statement is carried by the author's own lemma environment in the chapter source.\par\medskip
\noindent\textit{Allowed standard prerequisites (editorial).}\begin{itemize}
\item rig-surjectivity of morphisms of formal algebraic spaces tested by adic morphisms from complete DVRs after extension
\item base change preserves local finite type (Formal Spaces, Lemma formal-spaces-lemma-base-change-finite-type)
\item the diagonal of a base change is the base change of the diagonal, preserving rig-etale and rig-surjective
\item the dichotomy for the flat locus of a morphism to Spf(R) with rig-etale rig-surjective diagonal: empty or everything (lemma-flat-and-diagonal-rig-surjective of the same chapter)
\end{itemize}\medskip
\ifdefined\hypertarget\hypertarget{sourcepointer-1}{}\fi
\noindent\textit{Source reference.} Original lemma 0AQ5: sources/stacks-chow/tags.txt, source lines 13498--13498. Prerequisite pointer invoked at proof line(s) 5931; it names a standard fact of a sibling chapter that the counted argument uses, and every step of the counted proof built on it is retained verbatim in the delivered body.\par\medskip
\ifdefined\hypertarget\hypertarget{sourcepointer-2}{}\fi
\noindent\textit{Source reference.} Original reference 0GD9: sources/stacks-restricted/pinned-original-source.tex, source lines 5832--5832. Prerequisite pointer invoked at proof line(s) 5934; it names a standard fact of the same chapter that the counted argument uses, and every step of the counted proof built on it is retained verbatim in the delivered body.\par\medskip
\begin{lemma}
\label{lemma-rig-monomorphism-rig-surjective}
Let $S$ be a scheme. Let $f : X \to Y$ be a morphism of formal algebraic
spaces. Assume
\begin{enumerate}
\item $X$ and $Y$ are locally Noetherian,
\item $f$ locally of finite type,
\item $\Delta_f : X \to X \times_Y X$ is rig-\'etale and rig-surjective.
\end{enumerate}
Then $f$ is rig surjective if and only if every adic morphism
$\text{Spf}(R) \to Y$ where $R$ is a complete discrete
valuation ring lifts to a morphism $\text{Spf}(R) \to X$.
\end{lemma}

\begin{proof}
One direction is trivial. For the other, suppose that $\text{Spf}(R) \to Y$
is an adic morphism such that there exists an extension of complete
discrete valuation rings $R \subset R'$ with
$\text{Spf}(R') \to \text{Spf}(R) \to X$ factoring through $Y$.
Consider the fibre product diagram
$$
\xymatrix{
\text{Spf}(R') \ar[r] \ar[rd] &
\text{Spf}(R) \times_Y X \ar[r] \ar[d]^p &
X \ar[d]^f \\
&
\text{Spf}(R) \ar[r] &
Y
}
$$
The morphism $p$ is locally of finite type as a base change of $f$, see
Formal Spaces, Lemma
\ref{formal-spaces-lemma-base-change-finite-type}.
The diagonal morphism $\Delta_p$ is the base change of
$\Delta_f$ and hence is rig-\'etale and rig-surjective.
By Lemma \ref{lemma-flat-and-diagonal-rig-surjective}
the flat locus of $\text{Spf}(R) \times_Y X$ over $R$
is either $\emptyset$ or equal to $\text{Spf}(R)$.
However, since $\text{Spf}(R')$ factors through it
we conclude it is not empty and hence
we get a morphism $\text{Spf}(R) \to \text{Spf}(R) \times_Y X \to X$
as desired.
\end{proof}

\end{document}
